
\FigureLayoutDeclare{img/2022/shanghai_spring/q11_solution_fig1.png}{8.0cm}{0bp 0bp 0bp 0bp}

\chapter{2022年上海卷（春）}

\section{填空题}

\begin{problem}
已知 \(z=2+i\)，则 \(\bar z=\fillinblank{}\).
\end{problem}

\begin{answer}
\(2-i\)
\end{answer}

\begin{solution}
由共轭复数的定义，\(\bar z=2-i\).
\end{solution}

\begin{problem}
已知 \(A=(-1,2)\)，\(B=(1,3)\)，则 \(A\cap B=\fillinblank{}\).
\end{problem}

\begin{answer}
\((1,2)\)
\end{answer}

\begin{solution}
由集合交集的定义，\(A\cap B=(1,2)\).
\end{solution}

\begin{problem}
不等式 \(\dfrac{x-1}{x}<0\) 的解集为\fillinblank{}.
\end{problem}

\begin{answer}
\(\{x\mid 0<x<1\}\)
\end{answer}

\begin{solution}
由分式不等式的运算性质，分子，分母异号，故解集为 \(\{x\mid 0<x<1\}\).
\end{solution}

\begin{problem}
已知 \(\tan\alpha=3\)，则 \(\tan\left(\alpha+\dfrac\pi4\right)=\fillinblank{}\).
\end{problem}

\begin{answer}
\(-2\)
\end{answer}

\begin{solution}
由和角公式，\(\tan\left(\alpha+\dfrac\pi4\right)=\dfrac{\tan\alpha+\tan\dfrac\pi4}{1-\tan\alpha\tan\dfrac\pi4}=\dfrac{3+1}{1-3}=-2\).
\end{solution}

\begin{problem}
已知方程组 \(\begin{cases}
x+my=2,\\
mx+16y=8
\end{cases}\) 有无穷解，则 \(m\) 的值为\fillinblank{}.
\end{problem}

\begin{answer}
\(4\)
\end{answer}

\begin{solution}
方程组有无穷解，故两条直线重合，系数对应成比例，从而 \(1\times16=m^2\)，且 \(1\times8=2m\)，解得 \(m=4\).
\end{solution}

\begin{problem}
已知函数 \(f(x)=x^3\) 的反函数为 \(y=f^{-1}(x)\)，则 \(f^{-1}(27)=\fillinblank{}\).
\end{problem}

\begin{answer}
\(3\)
\end{answer}

\begin{solution}
由 \(y=x^3\) 得 \(x=\sqrt[3]{y}\)，故 \(f^{-1}(x)=\sqrt[3]{x}\)，所以 \(f^{-1}(27)=3\).
\end{solution}

\begin{problem}
在 \(\left(x^3+\dfrac1x\right)^{12}\) 的展开式中，含 \(\dfrac1{x^4}\) 项的系数为\fillinblank{}.
\end{problem}

\begin{answer}
\(66\)
\end{answer}

\begin{solution}
展开式的通项为 \(T_{r+1}=\mathrm C_{12}^{r}(x^3)^{12-r}\left(\dfrac1x\right)^r=\mathrm C_{12}^{r}x^{36-4r}\). 令 \(36-4r=-4\)，得 \(r=10\)，所以所求系数为 \(\mathrm C_{12}^{10}=66\).
\end{solution}

\begin{problem}
在 \(\triangle ABC\) 中，\(\angle A=\dfrac\pi3\)，\(AB=2\)，\(AC=3\)，则 \(\triangle ABC\) 的外接圆半径为\fillinblank{}.
\end{problem}

\begin{answer}
\(\dfrac{\sqrt{21}}3\)
\end{answer}

\begin{solution}
由余弦定理，\(BC^2=AB^2+AC^2-2AB\cdot AC\cos A=4+9-2\times2\times3\times\dfrac12=7\)，得 \(BC=\sqrt7\). 由正弦定理，\(R=\dfrac{BC}{2\sin A}=\dfrac{\sqrt7}{2\sin\dfrac\pi3}=\dfrac{\sqrt{21}}3\).
\end{solution}

\begin{problem}
已知有 \(1\)，\(2\)，\(3\)，\(4\) 四个数字组成无重复数字，则比 \(2134\) 大的四位数的个数为\fillinblank{}.
\end{problem}

\begin{answer}
\(17\)
\end{answer}

\begin{solution}
千位为 \(3\) 或 \(4\) 时，有 \(2\mathrm A_3^3=12\) 个；千位为 \(2\)，百位为 \(3\) 或 \(4\) 时，有 \(2\mathrm A_2^2=4\) 个；千位为 \(2\)，百位为 \(1\) 时，只有 \(2143\) 比 \(2134\) 大，有 \(1\) 个. 所以共有 \(17\) 个.
\end{solution}

\begin{problem}
在 \(\triangle ABC\) 中，\(\angle C=\dfrac\pi2\)，\(AC=BC=2\)，\(M\) 为 \(AC\) 的中点，\(P\) 在线段 \(AB\) 上，则 \(\overrightarrow{MP}\cdot\overrightarrow{CP}\) 的最小值为\fillinblank{}.
\end{problem}

\begin{answer}
\(\dfrac78\)
\end{answer}

\begin{solution}
以线段 \(AB\) 的中点为原点，\(AB\) 所在直线为 \(x\) 轴，\(AB\) 的垂直平分线为 \(y\) 轴建立平面直角坐标系，则 \(M\left(\dfrac{\sqrt2}{2},\dfrac{\sqrt2}{2}\right)\)，\(C(0,\sqrt2)\). 设 \(P(x,0)\)，\(-\sqrt2\le x\le\sqrt2\)，则
\solutionfigure{\bitmapfigure[max width=0.50\linewidth]{img/2022/shanghai_spring/q10_solution_fig1.png}}

\(\overrightarrow{MP}\cdot\overrightarrow{CP}=\left(x-\dfrac{\sqrt2}{2},-\dfrac{\sqrt2}{2}\right)\cdot(x,-\sqrt2)=x^2-\dfrac{\sqrt2}{2}x+1\). 当 \(x=\dfrac{\sqrt2}{4}\) 时，取得最小值 \(\dfrac78\).
\end{solution}

\begin{problem}
已知双曲线 \(\dfrac{x^2}{a^2}-y^2=1\)，其中 \(a>0\)，双曲线右支上有任意两点 \(P_1(x_1,y_1)\)，\(P_2(x_2,y_2)\)，满足 \(x_1x_2-y_1y_2>0\) 恒成立，则 \(a\) 的取值范围是\fillinblank{}.
\end{problem}

\begin{answer}
\(a\ge1\)
\end{answer}

\begin{solution}
设 \(P_3(x_2,-y_2)\)，则 \(\overrightarrow{OP_1}\cdot\overrightarrow{OP_3}=x_1x_2-y_1y_2>0\)，所以 \(\overrightarrow{OP_1}=\overrightarrow{OP_3}\) 或 \(\angle P_1OP_3\) 为锐角. 设双曲线两条渐近线分别与第一，第四象限相交于 \(M\)，\(N\)，则
\solutionfigure{\bitmapfigure[max width=0.50\linewidth]{img/2022/shanghai_spring/q11_solution_fig1.png}}

\(0<\angle MON\le\dfrac\pi2\)，即 \(0<\dfrac1a\le\tan\dfrac\pi4=1\)，故 \(a\ge1\).
\end{solution}

\begin{problem}
已知 \(f(x)\) 为奇函数，当 \(x\in(0,1]\) 时，\(f(x)=\ln x\)，且 \(f(x)\) 关于直线 \(x=1\) 对称. 设 \(f(x)=x+1\) 的正数解依次为 \(x_1\)，\(x_2\)，\(x_3\)，\(\ldots\)，\(x_n\)，\(\ldots\)，则 \(\displaystyle\lim_{n\to\infty}(x_{n+1}-x_n)=\fillinblank{}\).
\end{problem}

\begin{answer}
\(2\)
\end{answer}

\begin{solution}
由 \(f(x)\) 为奇函数，得 \(f(x)=-f(-x)\) 且 \(f(0)=0\). 又由 \(f(x)\) 关于直线 \(x=1\) 对称，得 \(f(1+x)=f(1-x)\)，所以 \(f(2+x)=f(-x)=-f(x)\)，从而 \(f(4+x)=-f(2+x)=f(x)\)，故 \(f(x)\) 是以 \(4\) 为周期的周期函数. 作出 \(y=f(x)\) 与 \(y=x+1\) 的图象可知
\solutionfigure{\bitmapfigure[max width=0.50\linewidth]{img/2022/shanghai_spring/q12_solution_fig1.png}}

正数解相邻两项之差的极限为两条渐近线之间的距离 \(2\)，所以 \(\displaystyle\lim_{n\to\infty}(x_{n+1}-x_n)=2\).
\end{solution}

\section{单选题}

\begin{problem}
下列函数定义域为 \(\R\) 的是（\quad）

\choices
    {\(y=x^{-\frac12}\)}
    {\(y=x^{-1}\)}
    {\(y=x^{\frac13}\)}
    {\(y=x^{\frac12}\)}
\end{problem}

\begin{answer}
C
\end{answer}

\begin{solution}
\(x^{1/3}\) 在全体实数上有定义，其余函数的定义域均受限制，故选 C.
\end{solution}

\begin{problem}
已知 \(a>b>c>d\)，下列选项中正确的是（\quad）

\choices
    {\(a+d>b+c\)}
    {\(a+c>b+d\)}
    {\(ad>bc\)}
    {\(ac>bd\)}
\end{problem}

\begin{answer}
B
\end{answer}

\begin{solution}
由 \(a>c\)，\(b>d\)，两式相加得 \(a+c>b+d\)，故 B 正确. 例如 \(3>2>1>0\)，但 \(3+0=2+1\)，\(3\times0<2\times1\)，故 A，C 错误；例如 \(30>2>-1>-2\)，但 \(30\times(-1)<2\times(-2)\)，故 D 错误. 所以选 B.
\end{solution}

\begin{problem}
如图，上海海关大楼的上面可以看作一个正四棱柱，四个侧面有四个时钟，则相邻两个时钟的时针从 \(0\) 时转到 \(12\) 时（含 \(0\) 时不含 \(12\) 时）的过程中，能够相互垂直（\quad）次.

\bitmapfigure[max width=0.22\linewidth]{img/2022/shanghai_spring/q15_fig1.png}

\choices
    {\(0\)}
    {\(2\)}
    {\(4\)}
    {\(12\)}
\end{problem}

\begin{answer}
B
\end{answer}

\begin{solution}
根据正四棱柱相邻侧面的线线关系，时针在 \(3\) 时或 \(9\) 时相互垂直，在 \(0\) 时或 \(6\) 时相互平行，其余时刻不垂直，故共 \(2\) 次，选 B.
\end{solution}

\begin{problem}
已知等比数列 \(\{a_n\}\) 的前 \(n\) 项和为 \(S_n\)，前 \(n\) 项积为 \(T_n\)，则下列选项判断正确的是（\quad）

\choices
    {若 \(S_{2022}>S_{2021}\)，则数列 \(\{a_n\}\) 是递增数列}
    {若 \(T_{2022}>T_{2021}\)，则数列 \(\{a_n\}\) 是递增数列}
    {若数列 \(\{S_n\}\) 是递增数列，则 \(a_{2022}\ge a_{2021}\)}
    {若数列 \(\{T_n\}\) 是递增数列，则 \(a_{2022}\ge a_{2021}\)}
\end{problem}

\begin{answer}
D
\end{answer}

\begin{solution}
对于 A，取 \(a_1=-1\)，公比 \(q=-2\)，则 \(S_{2022}>S_{2021}\)，但数列不是递增数列；对于 B，取 \(a_1=-1\)，公比 \(q=2\)，则 \(T_{2022}>T_{2021}\)，但数列不是递增数列；对于 C，取 \(a_1=1\)，公比 \(q=\dfrac12\)，则 \(\{S_n\}\) 递增，但 \(a_{2022}<a_{2021}\). 对于 D，若 \(\{T_n\}\) 递增，则 \(T_n>T_{n-1}\)，从而 \(a_n>1\)，故 \(q\ge1\)，所以 \(a_{2022}\ge a_{2021}\)，选 D.
\end{solution}

\section{解答题}

\begin{problem}
如图，在圆柱 \(OO_1\) 中，底面半径为 \(1\)，\(AA_1\) 为圆柱 \(OO_1\) 的母线.

\begin{enumerate}
\item 若 \(AA_1=4\)，\(M\) 为 \(AA_1\) 的中点，求直线 \(MO_1\) 与底面的夹角大小；
\item 若圆柱的轴截面为正方形，求该圆柱的侧面积和体积.
\end{enumerate}

\bitmapfigure[max width=0.33\linewidth]{img/2022/shanghai_spring/q17_fig1.png}
\end{problem}

\begin{solution}
\begin{enumerate}
\item 直线 \(MO_1\) 与底面的夹角等于 \(\angle MO_1A_1\). 因为 \(AA_1\perp O_1A_1\)，且 \(MA_1=\dfrac12AA_1=2\)，\(O_1A_1=1\)，所以在直角三角形 \(MO_1A_1\) 中，\(\tan\angle MO_1A_1=\dfrac{MA_1}{O_1A_1}=2\)，故夹角为 \(\arctan2\).
\item 轴截面为正方形，故 \(AA_1=2\). 所以侧面积为 \(2\pi\times1\times2=4\pi\)，体积为 \(\pi\times1^2\times2=2\pi\).
\end{enumerate}
\end{solution}

\begin{problem}
已知数列 \(\{a_n\}\)，\(a_2=1\)，\(\{a_n\}\) 的前 \(n\) 项和为 \(S_n\).

\begin{enumerate}
\item 若 \(\{a_n\}\) 为等比数列，\(S_2=3\)，求 \(\displaystyle\lim_{n\to\infty}S_n\)；
\item 若 \(\{a_n\}\) 为等差数列，公差为 \(d\)，对任意 \(n\in\N^*\)，均满足 \(S_{2n}\ge n\)，求 \(d\) 的取值范围.
\end{enumerate}
\end{problem}

\begin{solution}
\begin{enumerate}
\item 由 \(S_2=a_1+a_2=3\) 得 \(a_1=2\)，所以等比数列的公比 \(q=\dfrac12\). 因此 \(S_n=\dfrac{2\left(1-\left(\frac12\right)^n\right)}{1-\frac12}=4\left(1-\left(\dfrac12\right)^n\right)\)，故 \(\displaystyle\lim_{n\to\infty}S_n=4\).
\item 由 \(a_2=1\) 得 \(a_1=1-d\)，且 \(S_{2n}=\dfrac{2n(a_1+a_{2n})}{2}=n(a_2+a_{2n-1})\ge n\)，所以 \((2n-3)d\ge-1\). 当 \(n=1\) 时，得 \(d\le1\)；当 \(n\ge2\) 时，得 \(d\ge\dfrac1{3-2n}\). 因为 \(\left\{\dfrac1{3-2n}\right\}\)（\(n\ge2\)）单调递增，且 \(-1\le\dfrac1{3-2n}<0\)，所以 \(d\ge0\). 综上，\(0\le d\le1\).
\end{enumerate}
\end{solution}

\begin{problem}
如图，矩形 \(ABCD\) 区域内，\(D\) 处有一棵古树，为保护古树，以 \(D\) 为圆心，\(DA\) 为半径划定圆 \(D\) 作为保护区域，已知 \(AB=30\mathrm{~m}\)，\(AD=15\mathrm{~m}\)，点 \(E\) 为 \(AB\) 上的动点，点 \(F\) 为 \(CD\) 上的动点，满足 \(EF\) 与圆 \(D\) 相切.

\begin{enumerate}
\item 若 \(\angle ADE=20^\circ\)，求 \(EF\) 的长；
\item 当点 \(E\) 在 \(AB\) 的什么位置时，梯形 \(FEBC\) 的面积有最大值，最大面积为多少？
\end{enumerate}

\bitmapfigure[max width=0.42\linewidth]{img/2022/shanghai_spring/q19_fig1.png}

（长度精确到 \(0.1\mathrm{~m}\)，面积精确到 \(0.01\mathrm{~m}^2\)）
\end{problem}

\begin{solution}
\begin{enumerate}
\item 设 \(EF\) 与圆 \(D\) 相切于 \(H\)，连接 \(DH\)，则 \(DH\perp EF\)，且 \(DH=AD=15\). 因为 \(AE=EH\)，所以 \(\angle ADE=\angle HDE=20^\circ\)，从而 \(EH=15\tan20^\circ\). 又 \(\angle HDF=90^\circ-2\angle ADE=50^\circ\)，所以 \(HF=15\tan50^\circ\). 因此 \(EF=EH+HF=15(\tan20^\circ+\tan50^\circ)\approx23.3\mathrm{~m}\).
\item 设 \(\angle ADE=\theta\)，则 \(AE=15\tan\theta\)，\(FH=15\tan(90^\circ-2\theta)\). 梯形 \(AEFD\) 的面积为
\(S=\dfrac{225}{4}\left(3\tan\theta+\dfrac1{\tan\theta}\right)\ge\dfrac{225\sqrt3}{2}\)，当 \(\tan\theta=\dfrac{\sqrt3}{3}\) 时取等号，此时 \(AE=15\tan\theta=5\sqrt3\approx8.7\mathrm{~m}\). 故当 \(AE=8.7\mathrm{~m}\) 时，梯形 \(FEBC\) 的面积最大，最大值为 \(15\times30-\dfrac{225\sqrt3}{2}\approx255.14\mathrm{~m}^2\).
\end{enumerate}
\end{solution}

\begin{problem}
在椭圆 \(\Gamma:\dfrac{x^2}{a^2}+y^2=1\) 中，直线 \(l:x=a\) 上有两点 \(C\)，\(D\)（\(C\) 点在第一象限），左顶点为 \(A\)，下顶点为 \(B\)，右焦点为 \(F\).

\begin{enumerate}
\item 若 \(\angle AFB=\dfrac\pi6\)，求椭圆 \(\Gamma\) 的标准方程；
\item 若点 \(C\) 的纵坐标为 \(2\)，点 \(D\) 的纵坐标为 \(1\)，则 \(BC\) 与 \(AD\) 的交点是否在椭圆上？请说明理由；
\item 已知直线 \(BC\) 与椭圆 \(\Gamma\) 相交于点 \(P\)，直线 \(AD\) 与椭圆 \(\Gamma\) 相交于点 \(Q\)，若 \(P\) 与 \(Q\) 关于原点对称，求 \(\lvert CD\rvert\) 的最小值.
\end{enumerate}
\end{problem}

\begin{solution}
\begin{enumerate}
\item 设椭圆焦距为 \(c\)，则 \(A(-a,0)\)，\(B(0,-1)\)，\(F(c,0)\). 由 \(\tan\angle AFB=\dfrac1c=\tan\dfrac\pi6=\dfrac{\sqrt3}{3}\)，得 \(c=\sqrt3\). 又 \(a^2=c^2+1=4\)，故椭圆 \(\Gamma\) 的标准方程为 \(\dfrac{x^2}{4}+y^2=1\).
\item 由 \(B(0,-1)\)，\(C(a,2)\) 得 \(BC:y=\dfrac3a x+1\)；由 \(A(-a,0)\)，\(D(a,1)\) 得 \(AD:y=\dfrac1{2a}(x+a)\). 联立得交点为 \(\left(\dfrac{3a}{5},\dfrac45\right)\). 代入椭圆方程左边，得 \(\dfrac{\left(\frac{3a}{5}\right)^2}{a^2}+\left(\dfrac45\right)^2=\dfrac9{25}+\dfrac{16}{25}=1\)，故交点在椭圆上.
\item 设 \(P(a\cos\theta,\sin\theta)\)，\(Q(-a\cos\theta,-\sin\theta)\). 由直线 \(BP\)，\(AQ\) 的方程，得 \(C\left(a,\dfrac{\sin\theta+1}{\cos\theta}-1\right)\)，\(D\left(a,\dfrac{2\sin\theta}{\cos\theta-1}\right)\). 令 \(t=\tan\dfrac\theta2\)，则
\(\lvert CD\rvert=2\left\lvert\dfrac1{1-t}+\dfrac1t-1\right\rvert\). 在相应取值范围 \(0<t<1\) 内，由基本不等式，\(\dfrac1{1-t}+\dfrac1t\ge4\)，所以 \(\lvert CD\rvert\ge6\)，故最小值为 \(6\).
\end{enumerate}
\end{solution}

\begin{problem}
已知函数 \(f(x)\)，甲变化：\(f(x)-f(x-t)\)；乙变化：\(\lvert f(x+t)-f(x)\rvert\)，\(t>0\).

\begin{enumerate}
\item 若 \(t=1\)，\(f(x)=2^x\)，\(f(x)\) 经甲变化得到 \(g(x)\)，求方程 \(g(x)=2\) 的解；
\item 若 \(f(x)=x^2\)，\(f(x)\) 经乙变化得到 \(h(x)\)，求不等式 \(h(x)\le f(x)\) 的解集；
\item 若 \(f(x)\) 在 \((-infty,0)\) 上单调递增，将 \(f(x)\) 先进行甲变化得到 \(u(x)\)，再将 \(u(x)\) 进行乙变化得到 \(h_1(x)\)；将 \(f(x)\) 先进行乙变化得到 \(v(x)\)，再将 \(v(x)\) 进行甲变化得到 \(h_2(x)\)，若对任意 \(t>0\)，总存在 \(h_1(x)=h_2(x)\) 成立，求证：\(f(x)\) 在 \(\R\) 上单调递增.
\end{enumerate}
\end{problem}

\begin{solution}
\begin{enumerate}
\item 由甲变化，\(g(x)=f(x)-f(x-1)=2^x-2^{x-1}=2^{x-1}\). 令 \(g(x)=2\)，解得 \(x=2\).
\item 由乙变化，\(h(x)=\lvert(x+t)^2-x^2\rvert=t\lvert2x+t\rvert\). 当 \(2x+t<0\)，即 \(x<-\dfrac t2\) 时，\(t\lvert2x+t\rvert\le x^2\) 恒成立；当 \(2x+t\ge0\)，即 \(x\ge-\dfrac t2\) 时，解 \(t(2x+t)\le x^2\)，得 \(x\le(1-\sqrt2)t\) 或 \(x\ge(1+\sqrt2)t\). 所以解集为 \(\left(-\infty,(1-\sqrt2)t\right]\cup\left[(1+\sqrt2)t,+\infty\right)\).
\item 记 \(u(x)=f(x)-f(x-t)\)，则 \(h_1(x)=\lvert u(x+t)-u(x)\rvert\)；记 \(v(x)=\lvert f(x+t)-f(x)\rvert\)，则 \(h_2(x)=v(x)-v(x-t)\). 题设的“对任意 \(t>0\)，总存在 \(x\)”并不能推出对任意 \(x\) 的单调性。反例为
\[
f(x)=\begin{cases}x,&x<0,\\-1,&x\ge0.\end{cases}
\]
它在 \((-\infty,0)\) 上单调递增；对任意 \(t>0\)，取 \(x<-t\)，则两次增量均为 \(t\)，从而 \(h_1(x)=h_2(x)=0\)，但 \(f\) 在 \(0\) 处向下跳跃，并非在 \(\R\) 上单调递增。因此原命题按字面不成立，不能给出所要求的证明.
\end{enumerate}
\end{solution}
